\makeatletter
\def\input@path{{./}{../}{../../}{../../../}{../../../../}{preamble/}{../preamble/}{../../preamble/}{../../../preamble/}{../../../../book/preamble/}}
\makeatother
\input{preamble_exams.tex}
\input{graphs/highlights.tex}
% \GTenHighlightsfalse % Uncomment to disable instructional highlights.
% Population/model, sample evidence, design threshold, and result roles.
\newcommand{\Pop}[1]{\HighlightStroke{GTenHLParameter}{#1}}
\newcommand{\Obs}[1]{\HighlightStroke{GTenHLValue}{#1}}
\newcommand{\Design}[1]{\HighlightStroke{GTenHLThird}{#1}}
\newcommand{\Result}[1]{\HighlightStroke{GTenHLOutcome}{#1}}

\title{Grade 11 Global Economics\\Term 3 --- Lesson 1\\Home Practice C6: Tests for a population proportion}
\author{}\date{}
\hypersetup{pdftitle={Grade 11 Global Economics - Term 3 - Lesson 1 - C6},pdfsubject={Tests for a population proportion}}
\begin{document}
\maketitle
\setcounter{tocdepth}{1}
\setcounter{secnumdepth}{2}
\phantomsection\label{toc}
\tableofcontents
\section*{Evaluated criteria}
\begin{description}
\item[\textbf{C6}] Conducts hypothesis tests for a population proportion using the Z-test, interpreting the test statistic in contexts such as default rates, investment success, or market participation.
\end{description}
\textbf{How to practise.} Define success precisely. Use the null proportion in the standard error, and assess conditions before computing Z. Attempt each task before opening its solution link. Use the return arrows to revisit the problem, group, or contents. All datasets and scenarios in this activity are illustrative.
\clearpage
\section{Introduction}\label{sec:introduction}
\subsection*{Highlight legend}
\phantomsection\label{tab:highlight-legend}
\begin{tabularx}{\linewidth}{@{}>{\raggedright\arraybackslash}p{.22\linewidth}>{\raggedright\arraybackslash}p{.25\linewidth}>{\raggedright\arraybackslash}X@{}}
\toprule
\textbf{Symbol or notation}&\textbf{Concept}&\textbf{Brief meaning}\\\midrule
$\Pop{p},\Pop{p_0}$ & \Pop{Population model} & Proportion and benchmark. \\
$\Obs{x},\Obs{n},\Obs{\hat p}$ & \Obs{Evidence} & Count, size, observed rate. \\
$\Pop{SE_0}$ & \Pop{Null standard error} & Null sampling variability. \\
$\Obs{z},\Design{\alpha}$ & \Design{Testing quantities} & Evidence and threshold. \\
\bottomrule
\end{tabularx}
\subsection{Introductory example 1: Market participation}\label{ex:intro-one}
Among 250 independently sampled firms, 190 participate in a digital market: $\Obs{\hat p}=\Obs{0.76}$. The reference is $\Pop{p_0}=\Pop{0.72}$. To investigate an increase, write $H_0:p=0.72$, $H_a:p>0.72$ before considering the observed rate.
\subsection{Introductory example 2: Default risk}\label{ex:intro-two}
Among 100 independently sampled loans, 15 default in one year. The reference is 10\%, so $SE_0=0.03$ and $z=(0.15-0.10)/0.03$. Positive Z indicates an observed excess, not a probability and not the number of defaulted loans.
\section{Notation Reference and Practice Procedure}\label{sec:notation-and-procedure}
\subsection{Notation reference}\label{tab:notation-reference}
\begin{tabularx}{\linewidth}{@{}>{\raggedright\arraybackslash}p{.16\linewidth}>{\raggedright\arraybackslash}p{.18\linewidth}>{\raggedright\arraybackslash}p{.28\linewidth}>{\raggedright\arraybackslash}X@{}}
\toprule
\textbf{Symbol or notation}&\textbf{Concept}&\textbf{Definition}&\textbf{Meaning in the introductory examples}\\\midrule
$\Pop{p},\Pop{p_0}$ & \Pop{Population model} & $p$ is the true success proportion; $p_0$ is the null value, both between 0 and 1. & \textbullet\ \hyperref[ex:intro-one]{Example 1}: $p_0=0.72$ participation.\par\textbullet\ \hyperref[ex:intro-two]{Example 2}: $p_0=0.10$ defaults. \\
$\Obs{x},\Obs{n},\Obs{\hat p}$ & \Obs{Evidence} & $x$ successes among $n$ observations give $\hat p=x/n$. & \textbullet\ \hyperref[ex:intro-one]{Example 1}: $190/250=0.76$.\par\textbullet\ \hyperref[ex:intro-two]{Example 2}: $15/100=0.15$. \\
$\Pop{SE_0}$ & \Pop{Null standard error} & $SE_0=\sqrt{p_0(1-p_0)/n}$; do not substitute $\hat p$ here. & \textbullet\ \hyperref[ex:intro-two]{Example 2}: $\sqrt{0.10(0.90)/100}=0.03$. \\
$\Obs{z},\Design{\alpha}$ & \Design{Testing quantities} & $z=(\hat p-p_0)/SE_0$. Compare a p-value with $\alpha$, or $z$ with a critical value. & \textbullet\ \hyperref[ex:intro-two]{Example 2}: positive Z measures excess defaults in null standard errors. \\
\bottomrule
\end{tabularx}
\subsection{Practice procedure}
\begin{enumerate}
\item Define the population, period, and binary success. State $H_0:p=p_0$ and the claim-based alternative; choose $\alpha$ in advance.
\item Check random independent observations. For sampling without replacement require a sample no larger than about 10\% of the population.
\item Require expected null counts $np_0\ge10$ and $n(1-p_0)\ge10$ for this normal approximation.
\item Calculate $\hat p=x/n$, $SE_0=\sqrt{p_0(1-p_0)/n}$, and $z=(\hat p-p_0)/SE_0$.
\item Use the planned tail for the p-value and critical rule. Reject or fail to reject $H_0$; distinguish evidence from the economic size of the difference.
\end{enumerate}
\section{General Worked Problems}\label{sec:general-worked}
\SectionProblemsTableOfContents{\SectionProblemLink{prob:samples}{Same rate, different sample sizes} \\
\SectionProblemLink{prob:default}{A two-sided default review}}
\subsection{Same rate, different sample sizes}\label{prob:samples}
Two separate surveys sample firms independently from very large populations using simple random sampling. Survey A finds 190 market participants among 250 firms; B finds 380 among 500. Both test an increase above 72\% at $\alpha=0.05$. Supplied: right critical 1.645; $P(Z\ge1.41)\approx0.0793$, $P(Z\ge1.99)\approx0.0233$.
\QuestionsTasks[sec:general-worked]{{6}/{sol:samples}/{Conduct both tests and explain why the same observed proportion can lead to different decisions.}}
\subsubsection*{Worked solution}\label{sol:samples}
For each target population, $p$ is its participation proportion; $H_0:p=0.72$, $H_a:p>0.72$. Independent random sampling from large populations is given. Null counts are 180 and 70 in A, 360 and 140 in B, all at least 10. Both have $\Obs{\hat p}=\Obs{0.76}$. Using the null rate,
\[
z_A=\frac{0.76-0.72}{\sqrt{0.72(0.28)/250}}\approx1.41,\qquad
z_B=\frac{0.76-0.72}{\sqrt{0.72(0.28)/500}}\approx1.99.
\]
A: $1.41<1.645$ and p-value about $0.079>0.05$, so \Result{fail to reject $H_0$}. B: $1.99>1.645$ and p-value about $0.023<0.05$, so \Result{reject $H_0$}. The larger sample has lower null standard error. Both estimate a 4 percentage-point increase; its economic size is unchanged. These are separate fixed samples, not permission to extend A until rejection occurs.
\secInicio[prob:samples][sec:general-worked]
\subsection{A two-sided default review}\label{prob:default}
A simple random sample of 100 loans from 10,000 loans finds 15 one-year defaults. Loan outcomes are assumed independent. The preplanned question is any difference from 10\% at $\alpha=0.05$. Supplied: two-tail critical values $\pm1.960$; $P(Z\ge1.67)\approx0.0475$.
\QuestionsTasks[sec:general-worked]{{6}/{sol:default}/{Conduct the test and explain how the two-sided p-value is obtained.}}
\subsubsection*{Worked solution}\label{sol:default}
Let $\Pop{p}$ be the population one-year default proportion. $H_0:p=0.10$, $H_a:p\ne0.10$. Sampling is random and $100<0.10(10{,}000)$; independence is assumed. Expected null counts are $100(0.10)=10$ and $100(0.90)=90$. Thus
\[
\hat p=15/100=0.15,\quad SE_0=\sqrt{0.10(0.90)/100}=0.03,\quad z=0.05/0.03\approx1.67.
\]
The two-sided p-value is approximately $2(0.0475)=0.095$, using both tails. Since $|1.67|<1.960$ and $0.095>0.05$, \Result{fail to reject $H_0$}. Evidence is insufficient for a population difference. Observed defaults are 5 percentage points above the reference, which may still motivate risk monitoring. The statistic measures a departure of about 1.67 null standard errors, not a 1.67\% default rate.
\secInicio[prob:default][sec:general-worked]
\section{Advanced Problems}\label{sec:advanced}
\SectionProblemsTableOfContents{\SectionProblemLink{prob:limits}{When the Z approximation is unsuitable}}
\subsection{When the Z approximation is unsuitable}\label{prob:limits}
An investment desk samples 20 independent investments and records 3 successes. The benchmark success proportion is $p_0=0.05$, and the preregistered question is an increase at $\alpha=0.05$. A colleague proposes a Z-test using $\sqrt{\hat p(1-\hat p)/20}$. In a second proposal the desk samples 200 investments from just one strongly connected trading network.
\QuestionsTasks[sec:advanced]{{6}/{sol:limits}/{Assess both proposals and explain what additional method or design is needed.}}
\subsubsection*{Worked solution}\label{sol:limits}
For the first sample, $p$ is the population success proportion, $H_0:p=0.05$, $H_a:p>0.05$, and $\hat p=3/20=0.15$. Expected null counts are $20(0.05)=1$ and $20(0.95)=19$; the first is below 10. A normal Z approximation is unsuitable, regardless of the observed count. Also, a proportion Z-test uses $p_0$, not $\hat p$, in its standard error. An exact binomial test under independent Bernoulli sampling, or a larger planned sample, is needed; no Z decision is justified here. In the network proposal, a larger count does not cure dependence. Obtain a sampling design or dependence model appropriate to connected investments before reporting a test. Numerical formulas cannot compensate for invalid independence assumptions.
\secInicio[prob:limits][sec:advanced]
\end{document}

